\section{Discussion}\label{s:discussion}

This systematic review aimed to evaluate the efficacy, safety, and overall impact of angiogenic therapies in AMD prevention, focusing on visual outcomes. Key findings across various studies, such as those by \citet{kertes2019,kertes2020,wykoff2018}, and others, highlighted the effectiveness of different treatment strategies.

The \TreatAndExtend regimen with ranibizumab, as demonstrated by \citet{kertes2019,kertes2020}, emerged as a viable alternative to traditional monthly dosing, offering comparable visual improvements with reduced treatment frequency. This suggests a potential reduction in treatment burden without compromising efficacy. \citet{russo2018} and \citet{semeraro2015} presented innovative combination therapies, indicating enhanced outcomes and potentially fewer injections. The findings from \citet{nunes2019} and the \citet{catt2012} suggest flexibility in choosing between ranibizumab and bevacizumab, as both demonstrated similar efficacy and safety profiles.

\paragraph{Safety}\label{s:safety}
Safety remains a critical consideration, with studies indicating a generally favorable safety profile across various angiogenic therapies. However, long-term and multicenter studies are needed for a more comprehensive understanding of the risks and benefits.

\paragraph{Strengths and Limitations of the Review}\label{s:strengths-imitations}
A major strength of this review is its comprehensive analysis of diverse \RCTs, providing robust evidence across varied patient demographics, which enhances the generalizability of the findings. The populations included Australia \cite{schmidt2011}, Brazil \cite{nunes2019,schmidt2011}, Canada \cite{kertes2019,kertes2020}, Israel \cite{schmidt2011}, Italy \cite{russo2018,semeraro2015}, Turkey \cite{eldem2015,schmidt2011}, USA \cite{wykoff2018,catt2012,ho2014}, and more. The focus on RCTs ensures a high level of evidence quality. However, the review is limited by its emphasis on short-term outcomes and variability in study designs and patient populations. These limitations should be considered when interpreting the results and applying them to broader clinical practice.